\section{Experimental Setup}
\label{app:experimental_setup}

This section describes the Qwen3.5-4B online training procedure for
\methodabbr{}. Training budgets and experiment-specific changes are
reported in \cref{app:generalization_efficiency,app:beyond_frontier_learning,app:behavioral_induction,app:ablations}.

\subsection{Model, tasks, and agent environment}

We initialize the policy from Qwen3.5-4B
\citep{qwen2026qwen35fourb} and train on SWE-rebench-767, a fixed
collection of 767 software-engineering tasks. Each task is executed in
a repository-specific container. The
agent can read, write, and edit files, search the repository, and execute
shell commands. Its system instruction is:
\begin{quote}\small
You are a software-engineering agent working in a checked-out repository.
Use the available tools to inspect files, make focused edits, and run commands.
Solve the task, validate when practical, then give a brief final answer.
\end{quote}
The environment supplies the task description and evaluates the resulting
patch. We allow at most 100 agent turns, 8,192 generated tokens per turn,
and a context window of 131,072 tokens. Context compaction is disabled.
The solver samples at temperature 1.0 with top-$p=1.0$.

Task order is shuffled and the dataset is traversed cyclically, with
reshuffling at each epoch. Because generation and environment execution
are asynchronous, completion order also influences batch composition.

\subsection{Retrospection generation and supervision}
\label{app:retrospection_generation}

After each attempted solution, the same policy generates four independent
retrospections, conditioned on the same task, trajectory summary, patch,
and test feedback. The default prompt used in these experiments requests
a \emph{short paragraph without lists} that identifies a consequential
assumption or decision, relates it to supporting or contradicting evidence, gives a correction when
appropriate, and states concrete triggers for applying the lesson.
It explicitly distinguishes the correctness of individual decisions
from the overall verdict and asks the model to acknowledge insufficient
evidence. The complete prompt and two generated examples are provided in
\cref{app:retrospection_examples}.

We sample retrospections at temperature 0.9 with top-$p=1.0$ and an
8,192-token output limit. The short-paragraph instruction specifies this
prompt's requested answer format, not the general form of retrospection,
and does not exclude the model's preceding thinking from supervision.
Both failed and successful attempts contribute
targets, without reward weighting or a semantic-quality filter.

The prompt reports the raw test outcome separately from reward adjustments
and the reason for termination. An attempt that passes its tests but
exhausts the turn budget therefore retains a passing test verdict;
budget exhaustion is reported separately. If raw test feedback is
unavailable, the verdict is marked unknown rather than inferred from
the shaped reward.

The no-verdict ablation removes the entire post-attempt verdict and
test-feedback block from the retrospection prompt, rather than only
the pass/fail label. Observations recorded during the task attempt
remain available, including any test output already present in that
trajectory.
The faster-retrospection experiment instead retains the verdict and uses
it to choose the instruction (\cref{app:faster_rollouts}).

The retrospective context is constructed from a bounded rendering of
the interaction, rather than the complete solver context.
\Cref{tab:retrospection_context} lists the character budgets.
For fields exceeding a budget, we retain the first two-thirds and last
one-third, separated by a truncation marker. Thought, action, and
observation fields are shortened independently before the trajectory-level
budget is applied. The full prompt and target must also fit the
131,072-token training sequence limit.

\begin{table}[H]
\centering
\small
\caption{Character budgets for constructing the retrospection input.
Budgets exclude inserted truncation markers and structural headings.}
\label{tab:retrospection_context}
\begin{tabularx}{0.85\linewidth}{@{}Xr@{}}
\toprule
Input component & Retained characters \\
\midrule
Task description & 4,000 \\
Thought, action, or observation within each step & 600 each \\
Concatenated trajectory summary & 24,000 \\
Final patch & 2,000 \\
Test output & 3,000 \\
\bottomrule
\end{tabularx}
\end{table}

We discard generations that terminate at the output-length limit,
empty or failed generations, and prompt--target pairs that exceed the
training sequence limit. Preprocessing removes surrounding whitespace,
trailing termination markers, and a leading empty thinking block, while
preserving nonempty thinking and the final answer. An end-of-turn token
and newline are appended to each accepted target. All target tokens are
supervised, including thinking and the end-of-turn suffix; all prompt
tokens are masked. Rejected targets are not replaced with action-token
supervision.

To detect degenerate outputs, training stops after two consecutive batches
in which the median target length falls below 100 tokens or more than
10\% of targets contain at most 10 tokens. No minimum-length filter
is applied to individual retrospections.

\subsection{Online optimization}
\label{app:roft_loss_implementation}

Each nominal update uses 64 source attempts and four retrospections
per attempt, yielding up to
256 targets before rejection. We make one optimization pass through
the retained targets and do not replay earlier batches. Source groups
are consumed in completion order. We discard groups that began generation
with weights more than one update behind the current policy.
Serving weights are refreshed after each optimizer update, so a
multi-turn attempt and its subsequent retrospections need not use a
single frozen policy snapshot.

The objective in \cref{eq:ro} averages over all supervised
tokens in the retained batch.
We optimize all model parameters with AdamW, a constant learning rate
of $10^{-6}$, and no warmup. The data gradient is scaled by a constant
factor of four before norm clipping.
Relative token weighting is unchanged.
We use neither a reference-policy KL penalty nor an entropy bonus.
Optimizer settings and generation hyperparameters are summarized in
\cref{tab:roft_hyperparameters}.

\begin{table}[H]
\centering
\small
\caption{Reference online training and generation hyperparameters for
\methodabbr{}. Experiment-specific settings take precedence. Batch sizes
are nominal; targets rejected during preprocessing incur no loss.}
\label{tab:roft_hyperparameters}
\begin{tabularx}{\linewidth}{@{}p{0.48\linewidth}X@{}}
\toprule
Hyperparameter & Value \\
\midrule
Initial model & Qwen3.5-4B \\
Training tasks & 767 \\
Optimizer updates & Specified per experiment \\
Source attempts per update & 64 \\
Solver attempts per task selection & 1 \\
Retrospections per source & 4 \\
Target batch / microbatch size & 256 / 1 \\
Optimization passes per batch & 1; no replay \\
Optimizer & AdamW \\
Learning rate / schedule / warmup & $10^{-6}$ / constant / none \\
Adam $(\beta_1,\beta_2)$ / $\epsilon$ & $(0.9,0.98)$ / $10^{-8}$ \\
Weight decay / gradient-norm clip & 0.1 / 1.0 \\
Loss reduction & Mean over retained target tokens \\
Reference-KL / entropy coefficients & 0 / 0 \\
Solver temperature / top-$p$ & 1.0 / 1.0 \\
Retrospection temperature / top-$p$ & 0.9 / 1.0 \\
Top-$k$ restriction & None \\
Maximum solver turns & 100 \\
Maximum generated tokens per solver turn & 8,192 \\
Maximum generated tokens per retrospection & 8,192 \\
Solver context / total response limit & 131,072 / 131,072 tokens \\
Training sequence limit & 131,072 tokens \\
Context compaction & Disabled \\
Maximum policy lag at generation start & One update \\
Training precision / gradient accumulation & BF16 / FP32 \\
Attention / hidden dropout & 0 / 0 \\
\bottomrule
\end{tabularx}
\end{table}

\subsection{Compute and checkpointing}

We use eight NVIDIA B200 GPUs: four for training and four for
generating solution attempts and retrospections. Training uses
four-way context parallelism and accumulates gradients over the
complete update.

We save checkpoints every five successful optimizer updates and at the
final update. Evaluated checkpoints are reported in
\cref{app:generalization_efficiency,app:single_task_test_performance,app:ablations}.
Held-out evaluations are conducted separately from the online training
loop, with no stored retrospection or added reflection step in the
solver context.
